\newpage
\thispagestyle{empty}
\vspace*{0.35\textheight}
\begin{center}
    {\Huge\textbf{CAPÍTULO II}} \\[0.5cm]
    {\Huge\textbf{MARCO TEÓRICO}}
\end{center}

\newpage
\chapter{Marco Teórico}

\section{Procesamiento de Lenguaje Natural}

El procesamiento de lenguaje natural proporciona técnicas para representar, limpiar e interpretar texto no estructurado. En este proyecto, estas técnicas se aplican para transformar la información de las convocatorias públicas en una forma adecuada para búsqueda semántica.

\section{Embeddings}

Los embeddings son representaciones vectoriales densas que codifican relaciones semánticas entre fragmentos de texto. Su utilidad radica en permitir que consultas formuladas con términos distintos recuperen documentos conceptualmente cercanos.

\section{Bases vectoriales}

PostgreSQL con la extensión pgvector permite almacenar embeddings y realizar búsquedas por similitud dentro del mismo sistema de base de datos. Esta elección favorece una arquitectura mantenible, integrando información estructurada, filtros SQL y recuperación semántica en una sola plataforma.

\section{Retrieval-Augmented Generation}

El enfoque Retrieval-Augmented Generation combina recuperación de contexto relevante con un modelo de lenguaje que utiliza ese contexto para formular respuestas más precisas y fundamentadas.

\section{LangChain como capa de orquestación}

LangChain se utilizará como capa de integración para desacoplar la selección del modelo de lenguaje, la construcción del contexto y la lógica de recuperación, facilitando cambios de proveedor sin reestructurar completamente la aplicación.
